% Shared preamble for the repository diagrams. Each figure is compiled twice by build.sh: once as is
% (light theme) and once with \darktheme defined, which swaps the palette for a dark background.
\documentclass[tikz,border=4pt]{standalone}
\usepackage[T1]{fontenc}
\usepackage{lmodern}
\renewcommand{\familydefault}{\sfdefault}
\usepackage{amsmath,amssymb}
\usepackage{sansmath}
\sansmath
% Narrow boxes read better with whole words than with hyphenated ones.
\hyphenpenalty=10000
\exhyphenpenalty=10000
\usetikzlibrary{arrows.meta,positioning,calc,shapes.misc}

% Fill colors encode the role of a step: amber = test or rule, teal = forward simulation,
% purple = optimization or candidate generation, green = state or returned action, red = adverse
% outcome, blue = other processing.
\ifdefined\darktheme
  \definecolor{ink}{HTML}{E6EDF3}
  \definecolor{edge}{HTML}{8B949E}
  \definecolor{blue}{HTML}{1B2E45}
  \definecolor{teal}{HTML}{153C39}
  \definecolor{green}{HTML}{1C3823}
  \definecolor{amber}{HTML}{3F3313}
  \definecolor{red}{HTML}{442022}
  \definecolor{purple}{HTML}{30264A}
  \definecolor{grey}{HTML}{262B30}
  \definecolor{yes}{HTML}{56D364}
  \definecolor{no}{HTML}{F47067}
\else
  \definecolor{ink}{HTML}{17202A}
  \definecolor{edge}{HTML}{5D6670}
  \definecolor{blue}{HTML}{DDEBF7}
  \definecolor{teal}{HTML}{DDF2EF}
  \definecolor{green}{HTML}{E4F2E5}
  \definecolor{amber}{HTML}{FFF0CC}
  \definecolor{red}{HTML}{F8DDDD}
  \definecolor{purple}{HTML}{EAE2F5}
  \definecolor{grey}{HTML}{ECEFF1}
  \definecolor{yes}{HTML}{236B3B}
  \definecolor{no}{HTML}{A33636}
\fi

\tikzset{
  every picture/.style={color=ink, font=\small},
  box/.style={draw=edge, rounded corners=3pt, align=center, inner sep=4pt,
              text width=#1, minimum height=12mm, text=ink},
  box/.default=2.7cm,
  proc/.style={box=#1, fill=blue},   proc/.default=2.7cm,
  sim/.style={box=#1, fill=teal},    sim/.default=2.7cm,
  opt/.style={box=#1, fill=purple},  opt/.default=2.7cm,
  out/.style={box=#1, fill=green},   out/.default=2.7cm,
  bad/.style={box=#1, fill=red},     bad/.default=2.7cm,
  rule/.style={box=#1, fill=amber},  rule/.default=2.7cm,
  ask/.style={box=#1, fill=amber, chamfered rectangle, chamfered rectangle xsep=2pt,
              rounded corners=0pt, line width=0.6pt},
  ask/.default=2.7cm,
  arr/.style={-{Latex[length=2mm]}, semithick, draw=edge},
  darr/.style={arr, dashed},
  lab/.style={font=\scriptsize\bfseries, inner sep=1.5pt},
  ylab/.style={lab, text=yes}, nlab/.style={lab, text=no},
  note/.style={font=\scriptsize, text=ink, align=center},
}

\begin{document}
% Overview of the agent on its three time scales: inputs prepared offline, the revalidation done once
% by train(), and the four stages that act() consults at every step. The environment sits below the
% agent so that the observation and the action close one loop around act().
% Each time scale has its own light tint; boxes are left white so that the text stays legible.
\ifdefined\darktheme
  \definecolor{offband}{HTML}{151B23}\definecolor{offedge}{HTML}{8696AD}
  \definecolor{trband}{HTML}{211B0F}\definecolor{tredge}{HTML}{D4A72C}
  \definecolor{actband}{HTML}{0F1B2B}\definecolor{actedge}{HTML}{4C8EDA}
  \definecolor{simfill}{HTML}{0F2420}\definecolor{simedge}{HTML}{4FB3A0}
  \definecolor{envfill}{HTML}{151B23}\definecolor{envedge}{HTML}{8B949E}
\else
  \definecolor{offband}{HTML}{E6EBF1}\definecolor{offedge}{HTML}{50627A}
  \definecolor{trband}{HTML}{FBEBC6}\definecolor{tredge}{HTML}{8F6505}
  \definecolor{actband}{HTML}{DCE8F6}\definecolor{actedge}{HTML}{1F5596}
  \definecolor{simfill}{HTML}{D3EDE5}\definecolor{simedge}{HTML}{2A7364}
  \definecolor{envfill}{HTML}{E6EBF1}\definecolor{envedge}{HTML}{3F4A56}
\fi
\tikzset{
  abox/.style={draw=#1, fill=paper, rounded corners=3pt, line width=\lwstage, align=center,
               inner sep=6pt, text width=3.7cm, text=ink},
  bandlab/.style={font=\footnotesize, anchor=north east, inner sep=4pt, text=ink},
  alab/.style={font=\scriptsize, text=ink, inner sep=2pt},
}
\newcommand{\astg}[1]{\tikz[baseline=(b.base)]\node[circle, fill=actedge, text=paper, inner sep=0pt,
  minimum size=4.2mm, font=\footnotesize\bfseries] (b) {#1};}
\begin{tikzpicture}
% bands
\fill[offband, rounded corners=5pt] (-2.45,6.15) rectangle (18.35,9.15);
\fill[trband, rounded corners=5pt] (-2.45,2.95) rectangle (18.35,5.65);
\fill[actband, rounded corners=5pt] (-2.45,-4.75) rectangle (18.35,2.45);
\node[bandlab] at (18.35,9.15) {\textcolor{offedge}{\textbf{Prepared offline}}, before the competition};
\node[bandlab] at (18.35,5.65) {\textcolor{tredge}{\textbf{\texttt{train()}}}, once, before the scenarios are scored};
\node[bandlab] at (18.35,2.45) {\textcolor{actedge}{\textbf{\texttt{act()}}}, at every five-minute step};
% offline inputs
\node[abox=offedge, minimum height=24mm] (rec) at (2.65,7.45) {\textbf{Posture record}\\[2pt]\footnotesize configurations of substation 79 that survived the outage of line 175};
\node[abox=offedge, minimum height=24mm] (wset) at (7.95,7.45) {\textbf{Withheld lines}\\[2pt]\footnotesize default\\$W=(106,131,33)$\\[1pt]alternative\\$W'=(162,110,61)$};
\node[abox=offedge, minimum height=24mm] (lib) at (13.25,7.45) {\textbf{Topology library}\\[2pt]\footnotesize 907 single-substation busbar assignments that relieved stressed states};
% train()
\node[abox=tredge, minimum height=15mm] (rev) at (7.95,4.15) {\textbf{Revalidation}\\[2pt]\footnotesize $W'$ replaces $W$ only if it scores significantly better on training scenarios};
% act() stages
\node[abox=actedge, minimum height=35mm] (s1) at (0,0.0) {\astg{1}\ \textbf{Scheduled-outage look-ahead}\\[3pt]\footnotesize if an announced outage would be severe, prepares the network before it starts};
\node[abox=actedge, minimum height=35mm] (s2) at (5.3,0.0) {\astg{2}\ \textbf{Preventive withholding}\\[3pt]\footnotesize keeps one attackable line per region out of service, so attacks there are canceled};
\node[abox=actedge, minimum height=35mm] (s3) at (10.6,0.0) {\astg{3}\ \textbf{Reactive control}\\[3pt]\footnotesize relieves loading by regime; when in danger, searches topology and dispatch actions};
\node[abox=actedge, minimum height=35mm] (s4) at (15.9,0.0) {\astg{4}\ \textbf{Extended-horizon safeguard}\\[3pt]\footnotesize checks the action of 3 over 8 steps; if it would end the episode, replaces it by an alternative that survives longer};
\node[draw=simedge, fill=simfill, rounded corners=3pt, line width=\lwstage, align=center, inner sep=6pt,
      text width=19.2cm, text=ink, minimum height=11mm] (ev) at (7.95,-2.95) {\textbf{Forward simulation} with \texttt{obs.simulate}: a candidate action is played for one step, then do-nothing,\\giving its peak loading $F_H$ and its cost $J_H$};
\node[font=\footnotesize\itshape, text=ink, anchor=south, inner sep=5pt, align=center] at (7.95,-4.75) {Stages 1 and 2 end the step when they act. Otherwise stage 3 chooses an action, and stage 4 checks it before it is sent.\\Dotted lines: the stage scores candidate actions by forward simulation.};
% environment, below the agent
\node[draw=envedge, fill=envfill, rounded corners=3pt, line width=\lwstage, align=center, inner sep=7pt,
      text width=6.4cm, text=ink] (env) at (7.95,-6.5) {\textbf{Grid2Op environment} \texttt{l2rpn\_idf\_2023}\\[2pt]\footnotesize 118 buses, three regions, adversarial line outages};
% arrows between the stages
\draw[arr] (s1.east)--node[alab,above]{no action}(s2.west);
\draw[arr] (s2.east)--node[alab,above]{no action}(s3.west);
\draw[arr] (s3.east)--node[alab,above]{action}(s4.west);
\foreach \s in {s1,s2,s3,s4}{\draw[draw=edge, dotted, thick] (\s.south)--(\s.south|-ev.north);}
% inputs feeding the stages
\draw[darr] (rec.south)--(2.65,2.7)--(0,2.7)--(s1.north);
\draw[arr] (wset.south)--(rev.north);
\draw[arr] (rev.south)--node[alab,right,pos=0.45]{$W$ used in play}(7.95,2.7)--(5.3,2.7)--(s2.north);
\draw[darr] (lib.south)--(13.25,2.7)--(10.6,2.7)--(s3.north);
% loop with the environment
\draw[arr] (env.west)--(-3.05,-6.5)--node[alab,above,rotate=90]{observation $o_t$}(-3.05,0.0)--(s1.west);
\draw[arr] (s4.east)--(18.95,0.0)--node[alab,above,rotate=-90]{action $a_t$}(18.95,-6.5)--(env.east);
\end{tikzpicture}
\end{document}
